\documentclass[11pt]{article}
\usepackage[T1]{fontenc}
\usepackage{lmodern,amsmath,amssymb,amsthm,booktabs,microtype}
\usepackage[letterpaper,margin=1.01in]{geometry}
\usepackage[hidelinks]{hyperref}
\usepackage{fancyhdr}
\hypersetup{pdftitle={Two fixed Mills minima and the zeros of Euler products},pdfauthor={Research note for Keio}}
\pagestyle{fancy}\fancyhf{}\lhead{Fixed shifted Mills minima}\rhead{Research note}\cfoot{\thepage}
\setlength{\headheight}{14pt}
\newtheorem{theorem}{Theorem}
\newtheorem{lemma}[theorem]{Lemma}
\newcommand{\PP}{\mathbb P}
\newcommand{\NN}{\mathbb N}
\newcommand{\Tr}{\operatorname{Tr}}
\newcommand{\floor}[1]{\lfloor #1\rfloor}
\newcommand{\Rea}{\operatorname{Re}}
\title{Two fixed Mills minima and the zeros of Euler products}
\author{Research note for Keio}
\date{November 3, 2026}
\begin{document}
\maketitle
\begin{abstract}
We revisit the least prime-representing constants for the unmodified integer exponents $3^k-1$ and $3^k+1$. Starting with Saito's Pisot reduction, we give short proofs of the transcendence statements in Keio's corpus manuscript 064. A local valuation argument handles the negative shift. For the positive shift, an archimedean bound and three original prime-floor tests suffice. We then separate two kinds of data. Both least constants and all primes in square-root neighborhoods of all their outputs survive a thinning whose Euler product has a prescribed off-line zero. In contrast, widths and order from the full branching set recover every sufficiently large prime exactly. The proposed implication from either fixed constant to the Riemann hypothesis remains unproved.
\end{abstract}
\section{The question for ordinary primes}
For a set $P\subseteq\PP$ and an exponent sequence $C$, put
\[
 W_P(C)=\{A>1:\floor{A^{C_k}}\in P\text{ for every }k\ge1\},\qquad
 \xi_P(C)=\min W_P(C)
\]
when the minimum exists. Fix $C_k^e=3^k+e$ for $e\in\{-1,1\}$, and put $\alpha_e=\xi_{\PP}(C^e)$. Saito's existence theorem applies to both schedules. His Corollary 2.7 under RH gives the transcendence of $\alpha_{-1}$ \cite{saito}. Keio's manuscript 064 states both transcendence results without RH \cite{keio}. Its affine return argument and transfer of local limit identities are the starting point here.
\begin{theorem}\label{th:arith}
Both $\alpha_{-1}$ and $\alpha_{+1}$ are transcendental.
\end{theorem}
Sections 2 and 3 rederive this result from published inputs. For either sign, given Theorem~\ref{th:arith}, proving the reverse implication $[\alpha_e\text{ transcendental}]\Rightarrow\mathrm{RH}$ is logically equivalent to proving RH. Section 4 tests the role of the ordinary prime system. Section 5 records how much more information the full branching set contains. We claim no priority for Theorem~\ref{th:arith} or for the general idea of prescribing Euler-product singularities.

\section{Cubic traces and a local identity}
Fix the sign $e$ and put $n_k=(3^k+e)/2$, $B=\alpha_e^2$. The map $A\mapsto A^2$ identifies the two admissible sets in increasing order, so $B=\xi_{\PP}((n_k))$. Assume $B$ is algebraic. The following specialization describes the only possible obstruction.
\begin{lemma}\label{lem:red}
$B$ is a cubic Pisot number, and $P_k=\Tr_{\mathbb Q(B)/\mathbb Q}(B^{n_k})=\floor{B^{n_k}}$ is a strictly increasing sequence of primes for all sufficiently large $k$.
\end{lemma}
\begin{proof}
Apply Proposition 3.1(ii)--(iv) of \cite{saito}, with its short-interval input $\theta=21/40$ (Theorem 1.6, based on \cite{bhp}). The first exponent is at least one. The successive ratios are at least $5/2$ and tend to $3>40/19$, so the index set of that proposition contains every sufficiently large $k$. Also $n_m\mid n_{3m}$ with quotient greater than one. An integer value of $B^{n_m}$ would therefore force a composite floor. The Pisot power $B^g$ supplied by the proposition has $g\mid n_k$ eventually. Since $n_{k+1}-3n_k=-e$, we have $g=1$. The degree lies between two and $1+40/17<4$; the quadratic case would require all large $n_k$ to be odd primes, although infinitely many are even. The trace formula follows from part (iv). The two small conjugates have modulus less than one, so $P_k=B^{n_k}+o(1)$ and eventually increases strictly.
\end{proof}
\begin{lemma}\label{lem:local}
Write $B_1,B_2,B_3$ for the conjugates in a finite splitting field over $\mathbb Q_3$. Let $U$ index the conjugates that are $3$-adic units. For some sign $a\in\{1,-1\}$ there are algebraic numbers $y_i$, $i\in U$, satisfying
\[
  \sum_{i\in U}y_i=a,\qquad y_i^{208}=B_i^{104e}.
\]
\end{lemma}
\begin{proof}
First we recover the return congruence of \cite{keio}. Fix a sufficiently large $k$ and put $p=P_k>3|N(B)|$. Let $f$ be the monic irreducible polynomial of $B$. The class of $X$ in $\mathbb F_p[X]/(f)$ is invertible. Write its order as $3^v d$, $3\nmid d$. Choose $h>0$ with $2d\mid 3^h-1$. Then $n_{k+h}-n_k=3^k(3^h-1)/2$. If $v\le k$, the two powers of $X$ coincide in this algebra. Taking the trace of the multiplication matrices gives $P_{k+h}\equiv P_k\pmod p$, contradicting the primality of $P_{k+h}>p$. The reduction to the product of residue fields has kernel a $p$-group on units. Hence $3^{k+1}\mid p^j-1$ for some $j\in\{1,2,3\}$. For $j=2$, at most one of $p-1,p+1$ is divisible by three. For $j=3$, $p\equiv1\pmod3$ and $v_3(p^2+p+1)=1$. Consequently $P_k\equiv\pm1\pmod {3^k}$.

Now take $k$ through multiples of 12. By compactness, after passing to a subsequence each $B_i^{n_k}$ has a limit and $P_k$ tends $3$-adically to a fixed sign $a$. Nonunits tend to zero. Each nonzero reduced root has degree at most three over $\mathbb F_3$, so its order divides $\operatorname{lcm}(2,8,26)=104$. Thus $B_i^{104}$ is a principal unit for $i\in U$ and $(B_i^{104})^{3^k}\to1$. This last assertion follows from the binomial expansion, which strictly increases the valuation of $u-1$ on cubing a principal unit in a finite extension. Since $208n_k=104(3^k+e)$, the limits give the displayed identities. Their power equations show that the $y_i$ are algebraic over $\mathbb Q$.
\end{proof}

\section{Eliminating the two obstructions}
\paragraph{Negative shift.}
Suppose $e=-1$. The identity in Lemma~\ref{lem:local} rules out $U=\varnothing$. If $U$ is proper, take a Galois automorphism over $\mathbb Q$ mapping some $B_i$, $i\in U$, to a nonunit conjugate. With $|U|=1$ precisely one term of the transformed sum has negative $3$-adic valuation. With $|U|=2$ the same is true because there is just one nonunit among the three conjugates. All other terms have valuation zero. The sum would have negative valuation, contradicting its value $a$. Thus all roots are units. Since $3^{12}\equiv1\pmod{208}$, for every $12\mid k$ their $n_k$-th powers each reduce to 1. Then $P_k\equiv0\pmod3$ for arbitrarily large $k$, contradicting Lemma~\ref{lem:red}.

\paragraph{Positive shift.}
Now let $e=1$. Pick one $i\in U$ and embed the algebraic numbers into $\mathbb C$ taking $B_i$ to the dominant real root $B$. The corresponding term has modulus $\sqrt B$ and each other term has modulus less than one. Lemma~\ref{lem:local} gives $\sqrt B<3$, hence $1<B<9$. The two small conjugates cannot both be real. At the infinitely many even $n_k$ their powers would have positive sum, making $P_k>B^{n_k}$ contrary to the floor equality. Write them as $u,\bar u$ and their minimal polynomial together with $B$ as
\[
 f(T)=T^3-aT^2+bT-c,\qquad a=B+2\Rea(u),\quad
 b=2B\Rea(u)+|u|^2,\quad c=B|u|^2.
\]
The possible integer coefficients satisfy $0\le a\le10$, $|b|\le18$, $1\le c\le8$. Also $\operatorname{disc}(f)<0$ and $f(c)<0<f(9)$, since $c<B$ and $f$ has exactly one real root. These necessary conditions leave 619 polynomials. The original first three floor conditions require $\floor{B^2},\floor{B^5},\floor{B^{14}}$ all to be prime. Appendix~A gives an exact integer certificate.
\begin{center}
\begin{tabular}{lrrrr}
\toprule
Test reached & coefficient bounds and signs & $B^2$ & $B^5$ & $B^{14}$\\
\midrule
Candidates remaining & 619 & 154 & 21 & 0\\
\bottomrule
\end{tabular}
\end{center}
This proves Theorem~\ref{th:arith}. The tests at $k=1,2,3$ use the original floor conditions; they require no trace equality at these early indices.

\paragraph{A smaller certificate.}
The fractional-part estimate in Proposition 3.1(i) of \cite{saito} yields a second exclusion with two explicit final integers. Put $R=|u|$ and $u=R e^{i\phi}$, $0<\phi<\pi$. For every $0<\kappa<17/40$, the ratio limit $3$ and $\theta=21/40$ give
\[
 |2R^{n_k}\cos(n_k\phi)|=|P_k-B^{n_k}|\ll_\kappa B^{-\kappa n_k}.
\]
Indeed $P_k\sim B^{n_k}$ and $(1-\theta)n_{k+1}/n_k-1\to17/40$. If $R>B^{-\kappa}$, then $\cos(n_k\phi)\to0$. The recurrence $n_{k+1}=3n_k-1$ would give $|\cos(n_{k+1}\phi)|\to|\sin\phi|>0$, a contradiction. Thus $R\le B^{-17/40}$. The positive integer $c=BR^2\le B^{3/20}<2$ must equal one. The archimedean bound now reads $\sqrt B\le1+2B^{-1/4}$, forcing $B<11/2$ (at $11/2$ the left side exceeds $7/3$ while the right side is smaller). Imposing $c=1$ and $f(11/2)>0$ on the same coefficient list yields
\[
 32\ \longrightarrow\ 11\ \longrightarrow\ 2\ \longrightarrow\ 0.
\]
The last two polynomials and their exact floors are
\begin{center}
\begin{tabular}{lrrr}
\toprule
$f$ & $\floor{B^2}$ & $\floor{B^5}$ & $\floor{B^{14}}$\\
\midrule
$T^3-3T^2-T-1$ & 11 & 443 & 25\,714\,875\\
$T^3-3T^2+2T-1$ & 5 & 67 & 134\,643\\
\bottomrule
\end{tabular}
\end{center}
Both last entries are multiples of three.

\section{Keeping the minima and prescribing a zero}
Let $S_0$ be the union of the two prime output sequences. Each has $O(\log\log(3X))$ terms up to $X$, since $C_k^e\asymp3^k$. Protect every prime in
\[
 S=\PP\cap\bigcup_{v\in S_0}[v-\sqrt v\log^2(2v),\ v+\sqrt v\log^2(2v)] .
\]
Then $S(X)=O(\sqrt X\log^2(2X)\log\log(3X))$. Indeed, a sufficiently large center whose interval meets $[1,X]$ is at most $2X$. The Euler factors on $S$ alone give a holomorphic nonvanishing product on $\Rea s>1/2$ by absolute convergence. Removing just those factors from the ordinary zeta function therefore preserves every zero, with multiplicity, in that half-plane.

\begin{theorem}\label{th:cut}
Fix $1/2<\beta<1$ and $t>0$. There exists $P\subset\PP$ containing $S$, with $\xi_P(C^e)=\alpha_e$ for both signs, such that $\pi_P(x)=\pi(x)+O(x^\beta)$. The Euler product $Z_P(s)=\prod_{p\in P}(1-p^{-s})^{-1}$ continues meromorphically to $\Rea s>\beta/2$ and vanishes at $\beta+it$.
\end{theorem}
\begin{proof}
At every sufficiently large unprotected prime set $w_p=(2+2\cos(t\log p))p^{\beta-1}<1$, and set $w_p=0$ elsewhere. Put $W(x)=\sum_{p\le x}w_p$. Delete a prime exactly when $\floor{W(x)}$ jumps there. For the deleted set $D$ we have $D(x)=\floor{W(x)}=O(x^\beta)$. The retained set $P$ contains every output. Shrinking the permissible prime set cannot introduce a smaller admissible base, so both minima are preserved.

For $z=s+1-\beta$ we obtain
\[
 Z_P(s)=e^{H(s)}\,
 \frac{\zeta(s)}
 {\zeta(z)^2\,\zeta(z+it)\,\zeta(z-it)}
 \qquad(\Rea s>\beta/2),
\]
with $H$ holomorphic there. To verify this, start with logarithms for $\Rea s>1$. Partial summation extends $\sum_p(\mathbf1_D(p)-w_p)p^{-s}$ holomorphically to $\Rea s>0$ because its summatory coefficients are bounded. Restoring the weights on $S$ incurs a holomorphic error on $\Rea s>\beta-1/2$. The full weighted prime sum is the sum of the first prime terms of $2\log\zeta(z)+\log\zeta(z+it)+\log\zeta(z-it)$. Their higher powers converge absolutely for $\Rea z>1/2$. Higher powers for $D$ converge absolutely for $\Rea s>\beta/2$ by the count for $D$. As $\beta-1/2<\beta/2$, exponentiating proves the factorization and the stated continuation. At $s=\beta+it$, $\zeta(z-it)$ has its pole and the other denominator factors are finite and nonzero on $\Rea z=1$ \cite{dlmf}. The zero has multiplicity $1+\operatorname{ord}_{\beta+it}\zeta$.
\end{proof}
The construction uses only ordinary primes. On any $(x,x+h]$ containing $M$ ordinary primes, it deletes at most $4x^{\beta-1}M+1$, by the rounding rule. Thus even relative prime counts on intervals with $x,M\to\infty$ agree asymptotically. The new zero belongs to $Z_P$. For ordinary $\zeta$ the Euler factors run over the full set $\PP$. Prescribed singularities of subproducts have precedents, in particular Seip's Section 6 \cite{seip}.

\section{What the full branching family recovers}
For either sign, list the nonempty first-level branches in increasing order. A branch is specified by a sequence of prime floor values; its children are the nonempty extensions by one floor value, again in increasing order. Its width is the diameter of the set of bases satisfying the entire schedule with that prefix. A floor label $v$ at exponent $J$ has nominal cell $[v^{1/J},(v+1)^{1/J})$.

\begin{theorem}\label{th:geom}
Let a first-level branch have sufficiently large prime label $p$. Pick a median nonempty child, and a median nonempty child of that child, using either median when needed. Let $d$ be the width at the resulting third-level branch. Write $(D,E,F,G)=(2,8,26,80)$ for the negative shift and $(4,10,28,82)$ for the positive shift. Then
\[
             p=\floor{(F d)^{-D/(F-1)}}.
\]
Every sufficiently large prime occurs at the first level. Thus the order and widths in the full tree, with the numerical labels suppressed, recover all sufficiently large primes.
\end{theorem}
\begin{proof}
Call a cell full if its closure is strictly inside all prior nominal cells in its prefix. The first level has no such restrictions. Write $u$ for any successive exponent ratio. Always $5/2\le u\le4$. By the short-interval bound of Baker, Harman and Pintz \cite{bhp}, for every large $x$ there are primes within $O(x^\theta)$ on either side, $\theta=21/40$. In the next power coordinate a full cell with label $v$ provides a window of length $\asymp v^{u-1}$. Since $u-1-u\theta\ge3/16>0$, uniformly many disjoint prime cells fit strictly inside. Any such cell extends indefinitely by the same estimate; nested compact subcells give a common point, and strict containment at the next step enforces the required floors. This applies to every sufficiently large first-level prime.

Let $q,r$ be the labels after the two median choices. At least two nonempty children lie on either side of each median for large $p$. Only the two extreme unit cells can straddle a parent boundary. Thus both selected cells are full, and $q\ge p^{E/D}+1$, $q+1\le(p+1)^{E/D}-1$. By the mean value theorem, for a positive constant $c$,
\[
 p+c p^{1-E/D}\le r^{D/F}\le p+1-c p^{1-E/D}.
\]
The first and last fully contained, continuable children in the $r$-cell are within $O(r^{(G/F)\theta}+1)$ of the two ends in the fourth power coordinate. In the third power coordinate this costs $O(r^{-\delta})$ at either end, where $\delta=(1-\theta)G/F-1\in(0,1)$. The base-coordinate diameter consequently satisfies
\[
 Fd=r^{-(1-1/F)}(1+O(r^{-\delta})).
\]
Here the relative variation of the derivative on $[r,r+1]$ is $O(1/r)$. As $r\asymp p^{F/D}$, the base inside the floor of the theorem differs from $r^{D/F}$ by $O(p^{1-\delta F/D})$. For the two schedules $\delta F-E$ equals $4$ and $19/20$, respectively. The error is smaller than the margin on both sides, giving the claimed integer.
\end{proof}
If $N(X)$ counts the integers at most $X$ decoded by the rule, with branches where the rule is undefined omitted, $N(X)=\pi(X)+O(1)$. Consequently RH is equivalent to $N(X)=\operatorname{li}(X)+O_\epsilon(X^{1/2+\epsilon})$ for every $\epsilon>0$, by the classical prime-counting criterion \cite{lagarias}. Proving this error bound from the transcendence of either extreme base remains open in this note.

\appendix
\section{An integer verifier for the positive shift}
The following standard Python 3 code is complete. A dyadic bracket of width $2^{-100}$ is validated by integer comparisons at each floor. No floating-point or probable-prime operations enter. A value below two, or one with a proper divisor at most 243, fails its floor test. A surviving intermediate value need not be certified prime for this exclusion argument. There are no final survivors.
\begingroup\small
\begin{verbatim}
scale = 1 << 100

def bracket(a, b, c):
    def value(t, d):
        return t**3 - a*t*t*d + b*t*d*d - c*d**3
    disc = a*a*b*b - 4*b**3 - 4*a**3*c - 27*c*c + 18*a*b*c
    if disc >= 0 or not value(c, 1) < 0 < value(9, 1):
        return None
    low, high = c*scale, 9*scale
    while high - low > 1:
        mid = (low + high) // 2
        if value(mid, scale) < 0:
            low = mid
        else:
            high = mid
    assert value(low, scale) < 0 < value(high, scale)
    return low, high

cases = []
small = []
for a in range(11):
    for b in range(-18, 19):
        for c in range(1, 9):
            pair = bracket(a, b, c)
            if pair is not None:
                cases.append(pair)
                # 8*f(11/2) is an integer.
                if c == 1 and 1331 - 242*a + 44*b - 8*c > 0:
                    small.append(pair)

def screen(pool):
    counts = [len(pool)]
    last = []
    for n in (2, 5, 14):
        keep = []
        for low, high in pool:
            m = low**n // scale**n
            assert m == (high**n - 1) // scale**n
            if n == 14:
                last.append(m)
            if m >= 2 and all(m % j for j in range(2, min(244, m))):
                keep.append((low, high))
        pool = keep
        counts.append(len(pool))
    return counts, sorted(last)

assert screen(cases)[0] == [619, 154, 21, 0]
assert screen(small) == ([32, 11, 2, 0], [134643, 25714875])
\end{verbatim}
\endgroup
\begin{thebibliography}{9}
\bibitem{bhp} R. C. Baker, G. Harman and J. Pintz, The difference between consecutive primes, II. \emph{Proc. London Math. Soc.} 83 (2001), 532--562. \href{https://doi.org/10.1112/plms/83.3.532}{doi:10.1112/plms/83.3.532}.
\bibitem{keio} Keio, \emph{Ternary trace dynamics and the transcendence of shifted Mills constants}. Verifiable Mathematical Reasoning Corpus, manuscript 064, 25 October 2026, Theorems A and C. \href{https://github.com/QEEDDD/Verifiable-Mathematical-Reasoning-Corpus/tree/ecd87de3348dcb706a76bb12fde722d7a348a012}{Repository snapshot}.
\bibitem{lagarias} J. C. Lagarias, An elementary problem equivalent to the Riemann hypothesis. \emph{Amer. Math. Monthly} 109 (2002), 534--543. \href{https://arxiv.org/abs/math/0008177}{arXiv:math/0008177}, introduction.
\bibitem{dlmf} NIST Digital Library of Mathematical Functions, \href{https://dlmf.nist.gov/25.10}{\S25.10}, Zeros of the Riemann zeta function.
\bibitem{saito} K. Saito, Transcendency of variants of Mills' constant. \emph{Ramanujan J.} 70 (2026), article 69. Numbering here refers to \href{https://arxiv.org/html/2508.16068v3}{arXiv:2508.16068v3}. \href{https://doi.org/10.1007/s11139-026-01443-0}{doi:10.1007/s11139-026-01443-0}.
\bibitem{seip} K. Seip, Universality and distribution of zeros and poles of some zeta functions. \emph{J. Anal. Math.} (2020), \S6. \href{https://doi.org/10.1007/s11854-020-0126-3}{doi:10.1007/s11854-020-0126-3}.
\end{thebibliography}
\end{document}
